\documentclass{article}
\usepackage[T1]{fontenc}
\usepackage{lmodern}
\usepackage{graphicx}
\usepackage{amsmath,amssymb}
\usepackage[a4paper,margin=2cm]{geometry}
\usepackage[absolute,overlay]{textpos}
\setlength{\TPHorizModule}{1mm}
\setlength{\TPVertModule}{1mm}
\usepackage[indonesian]{babel}
\usepackage{hyperref}
\setlength{\emergencystretch}{2em}
% unit: Halaman 31

\begin{document}

% === Halaman 31 (python/native) ===

Pengurangan Dua Titik di dalam EC pada GF(p)

Misalkan P(xp, yp) dan Q(xq, yq).  

Pengurangan: P – Q  = P + (-Q), yang dalam hal ini –Q(xq, -yq (mod p)). 

31

\end{document}
